% ======================================================================
% "Sách đóng chỉ" (thread-bound book) post layout mockup — 花栏 border
%
% Engine: XeLaTeX (Vietnamese diacritics need Unicode + fontspec, same
% reasoning as essays/thien-tang.tex).
%
% Compile:  xelatex tay-du.tex   (twice — frame/rules stabilize on pass 2)
%
% Single page, straight into the article: masthead, headline, body with
% right-margin commentary column, author's red chop stamp at the close.
% ======================================================================
\documentclass[11pt,a4paper]{article}

\usepackage{fontspec}
\usepackage{polyglossia}
\setmainlanguage{vietnamese}

\usepackage[margin=2.6cm,top=2.4cm,bottom=2.6cm]{geometry}
\usepackage[table]{xcolor}
\usepackage{tikz}
\usetikzlibrary{calc}
\usepackage{eso-pic}
\usepackage{array}
\usepackage{tabularx}
\usepackage{microtype}
\usepackage{lettrine}
\usepackage{fancyhdr}

% ---- Palette ----
\definecolor{giay}{HTML}{E8E0CC}     % nền giấy dó
\definecolor{mucnho}{HTML}{1F1B16}   % chữ mực nho
\definecolor{sontrien}{HTML}{A32C1A} % son triện — dùng rất tiết chế
\definecolor{cham}{HTML}{2B3A4A}     % phụ trợ chàm

\pagecolor{giay}
\color{mucnho}

% ---- Typefaces ----
% Thân bài: serif đọc dài, ưu tiên phông có phủ dấu tiếng Việt đầy đủ.
\IfFontExistsTF{Linux Libertine O}{\setmainfont{Linux Libertine O}}{%
  \IfFontExistsTF{EB Garamond}{\setmainfont{EB Garamond}}{%
    \IfFontExistsTF{Cambria}{\setmainfont{Cambria}}{%
      \setmainfont{Latin Modern Roman}}}}

% Tiêu đề: display hơi mang nét thư pháp. (Segoe Print/Script thiếu dấu
% tiếng Việt — đã thử và loại; Constantia có nét bút mảnh-đậm tương phản,
% phủ đủ dấu, và khác hẳn phông thân bài để tạo phân tầng thị giác.)
\IfFontExistsTF{Constantia}{\newfontfamily\displayfont{Constantia}[BoldFont={Constantia Bold}]}{%
  \newfontfamily\displayfont{Cambria}[BoldFont={Cambria Bold}]}
\newcommand{\displaybold}{\displayfont\bfseries}

% Ngày tháng / số trang: mono.
\IfFontExistsTF{Consolas}{\setmonofont{Consolas}}{\setmonofont{Latin Modern Mono}}

\setlength{\parindent}{1.1em}
\setlength{\parskip}{0pt}

% ---- 花栏 (hoa lan) — flower-motif frame, redrawn on every page ----
\tikzset{
  flowerdot/.pic={
    \foreach \a in {0,72,...,288}{
      \fill[mucnho,opacity=0.5] (\a:0.055) circle (0.042);
    }
    \fill[sontrien,opacity=0.4] (0,0) circle (0.018);
  }
}
\newcommand{\PageFrame}{%
\begin{tikzpicture}[remember picture,overlay]
  \coordinate (NW) at ($(current page.north west)+(1.55cm,-1.55cm)$);
  \coordinate (NE) at ($(current page.north east)+(-1.55cm,-1.55cm)$);
  \coordinate (SW) at ($(current page.south west)+(1.55cm,1.55cm)$);
  \coordinate (SE) at ($(current page.south east)+(-1.55cm,1.55cm)$);
  \draw[mucnho,line width=1.0pt] (NW) rectangle (SE);
  \draw[mucnho,line width=0.35pt] ($(NW)+(0.11cm,-0.11cm)$) rectangle ($(SE)+(-0.11cm,0.11cm)$);
  \foreach \t in {0,...,11}{
    \pic at ($(NW)!{\t/11}!(NE)$) {flowerdot};
    \pic at ($(SW)!{\t/11}!(SE)$) {flowerdot};
  }
  \foreach \t in {0,...,16}{
    \pic at ($(NW)!{\t/16}!(SW)$) {flowerdot};
    \pic at ($(NE)!{\t/16}!(SE)$) {flowerdot};
  }
  \pic[scale=1.5] at (NW) {flowerdot};
  \pic[scale=1.5] at (NE) {flowerdot};
  \pic[scale=1.5] at (SW) {flowerdot};
  \pic[scale=1.5] at (SE) {flowerdot};
\end{tikzpicture}%
}
\AddToShipoutPictureBG*{\PageFrame}

% ---- Footer: page number in mono, sitting just inside the frame ----
\pagestyle{fancy}
\fancyhf{}
\renewcommand{\headrulewidth}{0pt}
\renewcommand{\footrulewidth}{0pt}
\fancyfoot[C]{\ttfamily\footnotesize\color{mucnho}--- \thepage\ ---}
\setlength{\footskip}{1.05cm}

% ---- Masthead: khung kẻ đôi trên/dưới tên báo ----
\newcommand{\PaperName}{THANH'S JOURNAL}
\newcommand{\Masthead}{%
\begin{center}
  \rule{\textwidth}{0.7pt}\\[1.2pt]\rule{\textwidth}{0.25pt}\\[7pt]
  {\displaybold\fontsize{38}{40}\selectfont \PaperName}\\[5pt]
  {\ttfamily\footnotesize SỐ 01 --- THÁNG 8, 2026}\\[7pt]
  \rule{\textwidth}{0.25pt}\\[1.2pt]\rule{\textwidth}{0.7pt}
\end{center}}

% ---- Chop stamp — con triện đỏ, xoay lệch, cuối bài ----
\newcommand{\ChopStamp}{%
\begin{flushright}
\begin{tikzpicture}
  \node[rotate=-7, rounded corners=2pt, fill=sontrien, minimum width=1.55cm,
        minimum height=1.55cm, inner sep=0pt] (seal) {};
  \node[rotate=-7] at (seal.center) {%
    \color{giay}\displaybold\fontsize{15}{15}\selectfont
    \begin{tabular}{c}T\\[-1pt]T\end{tabular}};
\end{tikzpicture}
\end{flushright}}

\begin{document}

% ==================== BÀI CHI TIẾT — thẳng vào nội dung ====================
\thispagestyle{fancy}
\Masthead
\vspace{10pt}

\begin{center}
  {\displaybold\LARGE Đường đi Tây Trúc thỉnh kinh}\\[4pt]
  \rule{0.3\textwidth}{0.3pt}\\[4pt]
  {\ttfamily\footnotesize Bởi Tiên Thành \quad\textperiodcentered\quad Hà Nội}
\end{center}
\vspace{12pt}

\renewcommand{\arraystretch}{1}
\begin{tabularx}{\textwidth}{@{}X !{\color{mucnho}\vrule width 0.35pt} p{0.235\textwidth}@{}}

\parbox[t]{\linewidth}{%
\lettrine[lines=3,lhang=0.02,loversize=0.02]{T}{ây Du Ký} không phải là
một cuộc hành trình để đến, mà là một cuộc hành trình để trở thành.
Tôn Ngộ Không rời khỏi Hoa Quả Sơn như một khối đá nứt ra từ ngông
cuồng, và mười vạn tám nghìn dặm sau, thứ trở về Ngũ Hành Sơn không
còn là con khỉ đại náo thiên cung nữa\textsuperscript{(1)}. Tám mươi
mốt kiếp nạn không phải là tám mươi mốt chướng ngại vật đặt ra để làm
khó nhân vật, mà là tám mươi mốt lần bản ngã bị bào mòn, cho đến khi
thứ còn lại đủ nhẹ để mang tên \emph{Đấu Chiến Thắng Phật}.

Đường Tăng, ngược lại, gần như không thay đổi. Ông đi từ đầu đến cuối
với cùng một lòng tin, cùng một sự yếu đuối, cùng một khuynh hướng tin
nhầm yêu tinh là dân lành\textsuperscript{(2)}. Có lẽ đó chính là vai
trò của ông trong câu chuyện: không phải để trưởng thành, mà để làm
điểm neo bất biến, một trục thẳng đứng để ba đồ đệ còn lại xoay quanh
mà tự đo lường chính mình.}

&

\vspace{2pt}
{\color{cham}\footnotesize
\textbf{(1)} Trong bản gốc, Ngô Thừa Ân dùng cả trăm chương để phân
biệt hai trạng thái tâm của Tôn Ngộ Không --- trước và sau vòng kim
cô --- điều dễ bị bỏ qua khi chỉ xem qua bản chuyển thể truyền hình.

\vspace{7pt}
\textbf{(2)} Chi tiết lặp lại này thường bị đọc như một điểm yếu kịch
bản, nhưng cũng có thể đọc như phép thử: lòng từ bi không phân biệt
đối tượng mới là thứ khiến ông xứng đáng đi trọn đường.}

\end{tabularx}

\vspace{16pt}
\ChopStamp

\vfill
\begin{center}
  \rule{\textwidth}{0.25pt}\\[3pt]
  {\ttfamily\footnotesize\color{cham} \PaperName{} \textperiodcentered{} in cho tác giả}
\end{center}

\end{document}
